% =============================================================================
%  1-Introduccion.tex — Proemio, declaraciones y definiciones
% =============================================================================

Acuerdo interpersonal que celebran, por una parte, \textbf{\ParteA} (en adelante, \enquote{Rosario}) y, por la otra, \textbf{\ParteB} (en adelante, \enquote{Galileo}), a quienes en conjunto se les denominará \enquote{Las Partes}, al tenor de las siguientes declaraciones, definiciones y cláusulas.

\section{Declaraciones}

\begin{enumerate}[label=\Roman*., leftmargin=*, itemsep=0.4em]
  \item \textbf{Origen del acuerdo.} Tras la lectura del Acuerdo 1, Rosario expresó dos inconformidades e incomodidades relacionadas entre sí:
  \begin{enumerate}[label=\alph*), itemsep=0.2em]
    \item Los chistes o bromas sobre la muerte que la involucren.
    \item El respeto a las personas y a su identidad.
  \end{enumerate}

  \item \textbf{Reconocimiento.} Galileo reconoce que en ocasiones tiende a no respetar la identidad sexual o ideológica de otras personas, en especial de la comunidad LGBTQ+ y sus asociados. Reconoce también que, por desconocimiento del tema, una opinión expresada con sinceridad o de manera asertiva puede resultar ofensiva sin que se perciba así en el momento, y que esto daña la relación con Rosario.

  \item \textbf{Solicitud de acompañamiento.} Galileo solicita que se le explique, ayude y apoye a entender qué es y qué no es ofensivo para ciertas comunidades o personas, así como las maneras adecuadas de dirigirse a ellas con respeto.

  \item \textbf{Aceptación sobre el humor.} Rosario expresó que el humor de Galileo puede ser difícil de entender y que, si se le incluye sin previo aviso, puede resultarle ofensivo. Galileo accede sin mayor problema a no incluir a Rosario en bromas sobre la muerte o el suicidio, y propone un procedimiento para los casos en que una broma no se entienda o pueda resultar ofensiva.

  \item \textbf{Bilateralidad.} Las Partes acuerdan que los derechos y obligaciones de este acuerdo aplican a ambas. Aunque el caso predominante es de Galileo hacia Rosario, esto no exime a Rosario de las mismas obligaciones hacia Galileo.

  \item \textbf{Fin común.} Las Partes reconocen que ambos puntos persiguen un mismo fin: que el trato entre ellas, y hacia las personas con quienes conviven, sea digno y respetuoso, y que las diferencias de conocimiento o de humor no dañen la relación.

  \item \textbf{Voluntad y alcance.} Las Partes celebran este acuerdo de manera libre, voluntaria y de buena fe, con el propósito de cuidar la relación y el bienestar emocional de ambas. Reconocen que no constituye un contrato con validez jurídica, sino un compromiso personal y moral entre ellas, complementario al Acuerdo 1.
\end{enumerate}

\subsection*{Definiciones}

Para efectos de este acuerdo se entenderá por:

\begin{description}[leftmargin=0pt, itemsep=0.3em]
  \item[Identidad:] la forma en que una persona se reconoce y se expresa, incluyendo, de manera enunciativa, mas no limitativa, su identidad sexual, de género, ideológica, cultural o religiosa.
  \item[Comunidad LGBTQ+:] el conjunto de personas lesbianas, gays, bisexuales, trans, queer y de otras orientaciones sexuales e identidades de género, así como sus asociados y aliados.
  \item[Expresión ofensiva:] comentario, broma, gesto o trato que descalifica, ridiculiza o invalida la identidad de una persona, con o sin intención de hacerlo.
  \item[Explicación anticipada:] explicación que una Parte da a la otra, antes de una convivencia y dentro de lo posible, sobre lo que puede resultar ofensivo para una persona o comunidad y sobre las maneras respetuosas de dirigirse a ella.
  \item[Broma:] comentario o expresión con intención humorística.
  \item[Límite personal:] lo que cada Parte está dispuesta a aceptar dentro de una broma o conversación sin sentirse ofendida o vulnerada. Solo la propia Parte puede definirlo.
  \item[Pausa de aclaración:] interrupción breve de una broma para preguntar si se trata de una broma, conforme a la Cláusula Quinta.
\end{description}
